% =====================================================================
% 06_prop3.tex — Proposition 3: Home wealth share / NFA revaluation (WP-C)
% =====================================================================
\section{Proposition 3: Home wealth share and NFA revaluation}
\label{sec:prop3}

This section proves Proposition~3 of KL: how a flight-to-safety shock
$\dev{\cy}_1>0$ redistributes financial wealth across countries on impact, and
how it revalues Home's net foreign assets. Unlike Propositions~1--2, here the
portfolios are \emph{not} identical: we allow Home to be levered long in capital,
to hold a nonzero Foreign-bond position, and the Home government to have issued
safe debt. The result is the realized wealth-share identity that the engine
already records as~\eqref{eq:wsh-realized} (their app.~eq.~(92)); the work of the
proof is to expose its three components and to combine them with the realized
returns of Proposition~\ref{prop:returns}.

\subsection{Statement}

\begin{proposition}[Home wealth share and NFA revaluation]
\label{prop:wealth}
Consider the simplified economy of Definition~\ref{def:simplified} with portfolios
initially close to the symmetric benchmark and with Home government safe debt
$b^g_{H,s}$ not too far from zero, and consider a flight-to-safety shock
$\dev{\cy}_1>0$ in period~$1$ with $\widetilde{\widetilde{\dev k}}_1=0$ on
impact. Then, to first order, Home's financial wealth share is revalued on impact
by
\begin{equation}
\dev{\wsh}_1 = \Bigl(\frac{\qk\,k}{\sav}-1\Bigr)\bigl(\rk_1-\rr_1\bigr)
+\frac{b_F}{\sav}\bigl(\fr{r}_1-\D\dev{\rer}_1-\rr_1\bigr)
-\disc\,\frac{\zs}{1+\zs}\,\frac{b^g_{H,s}}{\sav}\,\dev{\cy}_1, \label{eq:prop3-wsh}
\end{equation}
where $\qk k$, $b_F$, and $b^g_{H,s}$ are the steady-state Home positions in
capital, Foreign bonds, and government safe debt (as shares of Home saving
$\sav$), and the realized returns $\rk_1,\rr_1,\fr{r}_1,\D\dev{\rer}_1$ are those
of Proposition~\ref{prop:returns}. Home's net foreign assets $\widehat{\nfa}_1$
are revalued in the same direction.
\end{proposition}

\noindent (We transcribe the proposition to match KL's main-paper statement
(p.~1664). As in Proposition~\ref{prop:returns}, $\D\dev{\rer}_1=\dev{\rer}_1$
because $\dev{\rer}_0=0$ at the deterministic steady state.)

\subsection{Proof}

\paragraph{The identity is the engine's realized wealth-share linearization.}
Equation~\eqref{eq:prop3-wsh} is \emph{verbatim} the engine's realized
wealth-share reduced form~\eqref{eq:wsh-realized}, which is KL's app.~eq.~(92).
We reproduce here the derivation of that reduced form (their derivation of (92)),
so that the three coefficients in \eqref{eq:prop3-wsh} are exposed, and then sign
each component using Proposition~\ref{prop:returns}.

\subparagraph{Step 1: linearize the wealth-share law of motion.} Home's financial
wealth share evolves according to the accounting identity that next period's Home
wealth equals the gross return on this period's Home portfolio, divided by the
gross return on world wealth (the setup's wealth-share law of motion; their
app.~eq.~(51), in re-scaled form). Writing $\dev{\wsh}_1$ as the deviation of the
realized period-$1$ wealth share from its steady-state value and linearizing the
portfolio return around the symmetric steady state, the realized wealth share is
the steady-state-weighted sum of the \emph{excess} realized returns Home earns on
each asset it is differentially exposed to, plus the seigniorage Home earns on the
safe debt it has issued:
\begin{equation}
\dev{\wsh}_1 = \sum_{j}\frac{x_j}{\sav}\bigl(R^{j}_1 - R^{\text{ref}}_1\bigr)
\;+\;(\text{seigniorage on Home safe debt}), \label{eq:prop3-decomp-schema}
\end{equation}
where $x_j$ is the steady-state Home position in asset $j$, $R^{j}_1$ its realized
return, and $R^{\text{ref}}_1$ the return on the reference (safe dollar bond)
funding leg. The three differential exposures in the simplified economy are
(i) leverage in capital, (ii) the Foreign-bond position, and (iii) Home-issued
safe government debt; we take them in turn.
\gapflag{KL do not display the linearization of (51) into (92) term by term; they
write (92) directly as ``log-linearizing the realized evolution of Home's wealth
share implies.'' The schema~\eqref{eq:prop3-decomp-schema} is the standard
budget-constraint linearization for a wealth-share / NFA accounting identity (cf.\
the Devereux--Sutherland portfolio method used in the setup, Section~\ref{sec:setup}):
to first order, the change in a portfolio's relative value is the
steady-state-weighted sum of excess realized returns on the legs in which the
portfolio differs from the reference benchmark. The three coefficients
$\frac{\qk k}{\sav}-1$, $\frac{b_F}{\sav}$, and
$-\disc\frac{\zs}{1+\zs}\frac{b^g_{H,s}}{\sav}$ are the steady-state portfolio
weights that multiply each excess return, evaluated at the symmetric benchmark.
We reconstruct the assignment of each weight to its excess-return leg below; the
weights themselves are read from KL's eq.~(92).}

\subparagraph{Step 2: the leverage-in-capital term.} Home holds steady-state
capital $\qk k$ financed partly by issuing safe (dollar) claims. Its net long
position in capital, relative to a portfolio fully funded by the safe leg, is
$\qk k-\sav$ (its capital holdings net of its own wealth $\sav$); as a share of
$\sav$ this is $\frac{\qk k}{\sav}-1$. The realized excess return on this levered
position is the realized capital return over the dollar-bond funding return,
$\rk_1-\rr_1$. Hence the first term of \eqref{eq:prop3-wsh},
\begin{equation}
\Bigl(\frac{\qk k}{\sav}-1\Bigr)\bigl(\rk_1-\rr_1\bigr). \label{eq:prop3-T1}
\end{equation}
Sign on impact: with Home levered long capital, $\frac{\qk k}{\sav}-1>0$; and from
Proposition~\ref{prop:returns}(a)--(b), the realized excess return on capital is
negative, $\rk_1-\rr_1<0$ (zero capital return minus a positive dollar-bond
return under flexible wages, or a negative capital return minus a positive
dollar-bond return under predetermined wages). Therefore
\eqref{eq:prop3-T1}$<0$: Home \emph{loses} wealth share because the risky capital
it is levered into underperforms the dollar bonds it is short.

\subparagraph{Step 3: the Foreign-bond term.} Home holds steady-state Foreign
bonds $b_F$ (financed at the dollar-bond rate). The realized excess return on this
position, in Home consumption units, is the Home-currency Foreign-bond return over
the dollar-bond return, $\fr{r}_1-\D\dev{\rer}_1-\rr_1$ (the $\D\dev{\rer}_1$
converts the Foreign return into Home units). As a share of $\sav$ the position is
$\frac{b_F}{\sav}$, giving the second term of \eqref{eq:prop3-wsh},
\begin{equation}
\frac{b_F}{\sav}\bigl(\fr{r}_1-\D\dev{\rer}_1-\rr_1\bigr). \label{eq:prop3-T2}
\end{equation}
Sign on impact: from Proposition~\ref{prop:returns}(c),
$\fr{r}_1-\D\dev{\rer}_1-\rr_1<0$ (Foreign bonds underperform dollar bonds, both
because of the dollar's deflation and, under nominal rigidity, the real dollar
appreciation). With $b_F>0$ (Home long Foreign bonds), \eqref{eq:prop3-T2}$<0$:
Home loses wealth on its Foreign-bond position as well.

\subparagraph{Step 4: the seigniorage term.} The Home government has issued safe
dollar debt $b^g_{H,s}$ (a \emph{negative} position from Home's private-sector
perspective: $b^g_{H,s}<0$ means Home is a net issuer of safe claims). A share
$\frac{\zs}{1+\zs}$ of this safe debt is held by the Rest of the World. When the
safety shock raises the convenience yield $\dev{\cy}_1$, the issuer of safe debt
earns \emph{seigniorage}: it borrows at a convenience-discounted rate, and the
discount widens by $\dev{\cy}_1$ to first order. The Home private sector captures
the seigniorage on the share held abroad, discounted by $\disc$ (the safe debt
pays off next period). This is the third term of \eqref{eq:prop3-wsh},
\begin{equation}
-\disc\,\frac{\zs}{1+\zs}\,\frac{b^g_{H,s}}{\sav}\,\dev{\cy}_1. \label{eq:prop3-T3}
\end{equation}
Sign on impact: when Home is a net issuer, $b^g_{H,s}<0$, so
$-\disc\frac{\zs}{1+\zs}\frac{b^g_{H,s}}{\sav}>0$, and with $\dev{\cy}_1>0$ this
term is \emph{positive}: Home \emph{gains} wealth share through the seigniorage it
earns on the safe debt it has issued to the Rest of the World.
\gapflag{The coefficient $-\disc\frac{\zs}{1+\zs}\frac{b^g_{H,s}}{\sav}$ is read
directly from KL's eq.~(92); we supply its interpretation. The factor
$\frac{\zs}{1+\zs}$ is the RoW population share (so it is the fraction of Home
safe debt held \emph{abroad}, where the convenience rent is a transfer to Home
rather than an internal wash); the discount $\disc$ converts the next-period
payoff to present value; and $\dev{\cy}_1$ is the first-order widening of the
convenience discount. The convenience term enters the household Euler equation as
$(1+r_{t+1})/(1-\cy_t)$ (Section~\ref{sec:engine}, discussion after
\eqref{eq:euler-r}), so a one-unit rise in $\dev{\cy}_1$ lowers the real cost of
the safe debt by one unit to first order; the issuer's private-sector wealth
gains $\disc\frac{\zs}{1+\zs}\frac{|b^g_{H,s}|}{\sav}\dev{\cy}_1$. We reconstruct
the sign and economic content; the magnitude is KL's.}

\subparagraph{Step 5: assemble.} Summing \eqref{eq:prop3-T1}, \eqref{eq:prop3-T2},
and \eqref{eq:prop3-T3} reproduces \eqref{eq:prop3-wsh}, i.e.\ the engine's
\eqref{eq:wsh-realized}:
\[
\dev{\wsh}_1 = \underbrace{\Bigl(\tfrac{\qk k}{\sav}-1\Bigr)\bigl(\rk_1-\rr_1\bigr)}_{<0}
+\underbrace{\tfrac{b_F}{\sav}\bigl(\fr{r}_1-\D\dev{\rer}_1-\rr_1\bigr)}_{<0}
\underbrace{-\disc\tfrac{\zs}{1+\zs}\tfrac{b^g_{H,s}}{\sav}\dev{\cy}_1}_{>0\text{ when }b^g_{H,s}<0},
\]
the realized wealth-share revaluation, with the first two components negative and
the third positive (for a net issuer) on impact of a positive safety shock.

\paragraph{The ``close to the symmetric benchmark'' qualifier.} The linear
decomposition \eqref{eq:prop3-wsh} is valid because the coefficients
$\frac{\qk k}{\sav}-1$, $\frac{b_F}{\sav}$, and $\frac{b^g_{H,s}}{\sav}$ are
\emph{steady-state} portfolio weights, evaluated at the symmetric benchmark. KL's
qualifier ``portfolios initially close to the symmetric benchmark and $b^g_{H,s}$
not too far from zero'' is exactly what licenses this first-order treatment: it
keeps the deviations $\dev{\wsh}_1$ small enough that the realized returns
$\rk_1,\rr_1,\fr{r}_1,\D\dev{\rer}_1$ --- which were computed in
Proposition~\ref{prop:returns} \emph{at} the symmetric benchmark
($\Eone\dev{\wsh}_2=0$) --- remain valid inputs to \eqref{eq:prop3-wsh} to first
order. Were the steady-state portfolio far from symmetric, the realized returns
would themselves depend on the wealth-share dynamics (through the
$\Eone\dev{\wsh}_2$ terms in the engine), and the clean three-way decomposition
would break.
\gapflag{KL phrase the qualifier as ``portfolios initially close to the symmetric
benchmark.'' We make explicit \emph{why} it is needed: it decouples the
realized-return inputs (computed at symmetry in Proposition~\ref{prop:returns})
from the wealth-share output, so that the portfolio weights enter only as
constant coefficients. This is the standard first-order portfolio logic; KL leave
it implicit.}

\paragraph{NFA revaluation.} Home's net foreign assets are revalued in the same
direction as the wealth share. The engine's NFA reduced form is
\eqref{eq:nfa-solved}, $\widehat{\nfa}_1=\sav[\Eone\dev{\wsh}_2-\frac{\zs}{1+\zs}\frac{1}{1-\cshare}\Eone\dev{\tot}_2]$,
which is the linearization of the NFA definition (their app.~eq.~(87),
\eqref{eq:nfa-raw}). The realized revaluation channel that drives $\dev{\wsh}_1$
in \eqref{eq:prop3-wsh} feeds NFA through the same realized excess returns: a
fall in Home's wealth share on impact is, to first order, a fall in the market
value of Home's net foreign asset position. Concretely, the realized
return-differential terms~\eqref{eq:prop3-T1}--\eqref{eq:prop3-T2} are a
\emph{valuation} loss on Home's external balance sheet (Home's external assets,
levered long risky and long Foreign bonds, lose value relative to its external
liabilities, the safe dollar claims it has issued), exactly the revaluation that
moves $\widehat{\nfa}_1$ down; the seigniorage term~\eqref{eq:prop3-T3} is a
partial offset. Hence ``Home's NFA are revalued in the same way.'' \qed
(\,Proposition~\ref{prop:wealth}\,)

\subsection{Discussion}

\paragraph{Two transfers out, one rent in.} Proposition~3 decomposes the wealth
redistribution in a flight to safety into three pieces. The first two ---
\eqref{eq:prop3-T1} and \eqref{eq:prop3-T2} --- are a direct consequence of
Proposition~\ref{prop:returns}: Home's wealth falls in its \emph{leveraged
capital} position and in its \emph{Foreign-bond} position, because both of these
risky assets underperform the safe dollar bonds that fund them exactly when the
safety shock hits. The third --- \eqref{eq:prop3-T3} --- runs the other way:
Home's wealth \emph{rises} in the safe debt issued by its government, because Home
earns the convenience rent (seigniorage) on the share of that safe debt held by
the Rest of the World. The net revaluation is the sum: on impact, the two
valuation losses dominate for a country that is levered long risky assets and
short safe debt, so Home loses wealth and transfers it to the Rest of the World.

\paragraph{Connection to the target facts: the US as the world's insurer.} This
is the model's account of the United States as the ``world's insurer'' and of the
valuation channel of US external adjustment. The US runs an external balance sheet
that is \emph{long risky, short safe}: it holds foreign equity and risky claims
(here, levered capital and Foreign bonds) financed by issuing safe dollar claims
(Treasuries and dollar bank liabilities) to the rest of the world. In a global
flight to safety, this balance sheet delivers a large \emph{valuation loss} ---
the US transfers wealth to the rest of the world precisely in bad global states,
which is the insurance it provides. Proposition~3 is the structural counterpart of
the empirical NFA-revaluation fact: US net foreign assets fall in value in risk-off
episodes, by an amount governed by the leverage weight $\frac{\qk k}{\sav}-1$ and
the Foreign-bond weight $\frac{b_F}{\sav}$, partially offset by the seigniorage
the US earns on its safe-debt issuance, $-\disc\frac{\zs}{1+\zs}\frac{b^g_{H,s}}{\sav}\dev{\cy}_1$.

\paragraph{The big\_cy reading.} A \emph{large absolute} convenience yield
$\dev{\cy}_1$ does two things to \eqref{eq:prop3-wsh}. Through
Proposition~\ref{prop:returns} it magnifies the realized excess-return terms
\eqref{eq:prop3-T1}--\eqref{eq:prop3-T2} (larger deflation and real appreciation
$\Rightarrow$ larger underperformance of risky assets), deepening the valuation
transfer the US makes as the world's insurer. But it also magnifies the
seigniorage term \eqref{eq:prop3-T3}: the larger the absolute convenience yield,
the more the US earns on its safe-debt issuance, partially insulating its wealth.
The balance between these is precisely the quantitative question big\_cy poses ---
whether a DiTella-magnitude convenience yield lets the model match the observed
US insurance role and home-currency bias \emph{without} the counterfactual US
risk-tolerance asymmetry that Kekre and Lenel otherwise require.
